\chapter{Growth and Development Accounting}\label{ch:b1:growth-accounting}

From 1980 to 2019 Singapore's output grew by $6.00\%$ a year. Measured with
the Penn World Table, $3.71$ points of that came from more capital, $1.42$
from more workers and $0.75$ from more schooling; only $0.12$ points --- $2.0\%$
of the growth --- were left for doing more with the same inputs.
% figdata: sources
Over the same years Germany's output grew by $1.72\%$ a year, and a third
of it came from productivity. % figdata: sources
\omterm{def:b1:growth-accounting:growth-accounting}{Growth accounting} splits growth into the parts due to each input and a
residual; \omterm{def:b1:growth-accounting:development-accounting}{development accounting} does the same for the gaps in income
between countries. Neither is a theory of growth. Both say where to look.

\begin{openingchart}{Average yearly growth of real GDP to 2019, split into
the contributions of capital, workers, human capital and productivity
(the residual), eight economies, from 1960 or from the first year the table
observes the labour share: Japan and Korea 1970, India 1975, Germany and
Singapore 1980, Brazil and China 1992. The diamond marks output growth;
Brazil's productivity term is negative.}{Penn World Table 11.0 (rgdpna, rnna, emp,
hc, labsh)}
\begin{tikzpicture}
\begin{axis}[omrecord, grid=none, xbar stacked, bar width=9pt, xmin=-0.5, xmax=8, xlabel={Per cent a year},
  symbolic y coords={Brazil, Germany, United States, Japan, India, Korea, Singapore, China}, ytick=data,
  y dir=reverse, width=0.72\linewidth, height=0.5\linewidth, enlarge y limits=0.08,
  legend style={font=\footnotesize, at={(0.5,-0.2)}, anchor=north, legend columns=5}, reverse legend=false]
  \addplot[fill=omDef, draw=none] table[y=country, x=capital, col sep=comma] {figdata/bachelor-1/21-growth-accounting/out/sources.csv};
  \addplot[fill=omThm, draw=none] table[y=country, x=labour, col sep=comma] {figdata/bachelor-1/21-growth-accounting/out/sources.csv};
  \addplot[fill=omRecord, draw=none] table[y=country, x=human, col sep=comma] {figdata/bachelor-1/21-growth-accounting/out/sources.csv};
  \addplot[fill=gray!45, draw=none] table[y=country, x=tfp, col sep=comma] {figdata/bachelor-1/21-growth-accounting/out/sources.csv};
  \legend{Capital, Workers, Human capital, Productivity}
\end{axis}
\begin{axis}[omrecord, grid=none, hide axis, xmin=-0.5, xmax=8,
  symbolic y coords={Brazil, Germany, United States, Japan, India, Korea, Singapore, China}, y dir=reverse,
  width=0.72\linewidth, height=0.5\linewidth, enlarge y limits=0.08]
  \addplot[only marks, mark=diamond*, mark size=2.6pt, black] table[y=country, x=output, col sep=comma] {figdata/bachelor-1/21-growth-accounting/out/sources.csv};
\end{axis}
\end{tikzpicture}
\end{openingchart}

Over their own periods, China grew fastest, at $8.93\%$ a year from 1992,
with $3.42$ points from productivity; Korea's $6.72\%$ from 1970 had $1.2$;
Brazil's $2.48\%$ from 1992 had $-0.79$: its output grew less than its
inputs. % figdata: sources
The periods differ, so the bars compare countries only roughly.

\begin{recall}
\omterm{def:b1:growth-accounting:growth-accounting}{Growth accounting} with logarithms and \omterm{def:b1:growth-accounting:tfp}{total factor productivity} were met
at school, as were institutions as a cause of prosperity. The Solow model
is \cref{ch:b1:solow-model}; elasticities of substitution and Cobb--Douglas
shares are in \cref{ch:b1:technology}.
\end{recall}

\section{The Solow residual}\label{sec:b1:growth-accounting:residual}

Let output be $Y = AK^{\alpha}(HL)^{1 - \alpha}$, where $H$ is human capital per worker. Taking
logarithms and differentiating with respect to time,
\[
g_Y = \alpha g_K + (1 - \alpha)(g_L + g_H) + g_A,
\]
where $g_X = \dot X/X$. With competitive factor markets, $\alpha$ is capital's share of
income (\cref{ch:b1:long-run-model}); $g_Y$, $g_K$, $g_L$ and $g_H$ are measured; $g_A$ is what
is left.

\begin{definition}[Growth accounting]\label{def:b1:growth-accounting:growth-accounting}
\index{growth accounting}
\emph{Growth accounting} splits the growth rate of output into the
contributions of each input --- its growth rate times its share of
income --- and a residual.
\end{definition}

\begin{definition}[Solow residual, total factor productivity]\label{def:b1:growth-accounting:tfp}
\index{Solow residual}\index{total factor productivity}
The \emph{Solow residual} is the part of output growth not accounted for by
the growth of measured inputs, $g_A = g_Y - \alpha g_K - (1 - \alpha)(g_L + g_H)$. Its level $A$ is
\emph{total factor productivity} (TFP): output per unit of combined inputs.
\end{definition}

In data the growth rates are yearly differences of logarithms, and the
shares change from year to year. (Before a country's first observed \omterm{def:b1:factor-markets:factor-share}{labour
share} the table repeats that value; the opening chart leaves those years
out, so each bar starts in a year with a measured share.) The \emph{Törnqvist} version weights each
year's growth by the average of the shares in the two adjacent years,
$\bar\alpha_t = \tfrac12(\alpha_t + \alpha_{t-1})$; the opening chart is built that way, with capital's
share equal to one minus the \omterm{def:b1:factor-markets:factor-share}{labour share} of \cref{ch:b1:factor-markets}.

The residual is everything the inputs miss: better technology, but also
better organisation, the reallocation of workers from farms to factories,
the quality of capital not captured by its measured quantity, and every
error in measuring output and inputs. It has been called a measure of our
ignorance, and the name is fair.

\begin{model}{Growth accounting}
\assumes A Cobb--Douglas aggregate technology with constant returns;
factors paid their \omterm{def:b1:technology:marginal-product}{marginal products}, so that income shares measure the
elasticities; inputs measured correctly.
\predicts The decomposition itself is an identity, not a causal statement
(a rise in $A$ also raises $K$ through investment). What it predicts is that
the residual is small where accumulation explains growth, and that it
shrinks as inputs are better measured (hours, schooling, the quality of
capital); and, if \omterm{def:b1:factor-markets:factor-share}{factor shares} measure the elasticities, that estimates
from quantities and from factor prices agree.
\tested Better measurement does shrink the residual of the fast growers of
East Asia (Young 1995); % ledger: b1:young1995
but estimates from factor prices exceed those from quantities by about 1
point a year in Taiwan and more than 2 in Singapore (Hsieh 2002; debate
below). % ledger: b1:hsieh2002
\breaks Shares are not constant; capital and its quality are hard to
measure; the residual absorbs all measurement errors.
\end{model}

\section{The sources of growth}\label{sec:b1:growth-accounting:sources}

The fast growers of East Asia grew mostly by accumulating. In Hong Kong,
Singapore, Korea and Taiwan, output per person grew by $6$ to $7\%$ a year for
two and a half decades, but output per effective worker outside
agriculture grew by only $3$ to $4\%$; once the rise of investment rates,
participation and schooling is counted, productivity growth was within
the range of other economies (Young 1995). % ledger: b1:young1995
The opening chart, over later periods, finds the same pattern for
Singapore, and a larger residual for Korea and China.

\begin{debate}[How much of East Asian growth was productivity?]
Primal estimates, from quantities, find little productivity growth in
Singapore. Dual estimates, from factor prices, look at the return to
capital: if capital had piled up as fast as the national accounts say, its
return should have fallen, and in Singapore it did not. Dual estimates
exceed the primal ones by about $1$ point a year in Taiwan and by more than $2$
points in Singapore, and are close for Korea and Hong Kong (Hsieh 2002).
% ledger: b1:hsieh2002
Each method has its weakness: the primal one depends on capital stocks
built from investment data, the dual one on measured returns and
assumptions about how they are set. The disagreement is about measurement,
and it is not settled.
\end{debate}

How the labour input is measured matters too. Counting hours rather than
persons, the American labour contribution falls from $0.88$ to $0.78$ points a
year and the residual rises from $0.74$ to $0.84$; in Germany since 1980,
where hours per worker fell, the labour contribution turns from $0.26$ to
$-0.16$ and the residual rises from $0.56$ to $0.98$. % figdata: hours

\section{Levels accounting}\label{sec:b1:growth-accounting:levels}

The same \omterm{def:b1:technology:production-function}{production function} can split the \emph{level} of income per
worker, $y = A k^{\alpha} h^{1 - \alpha}$, with $k = K/L$ and $h$ human capital per worker. The part
due to inputs is $y_{KH} = k^{\alpha}h^{1 - \alpha}$; the rest is $A$.

\begin{definition}[Development accounting]\label{def:b1:growth-accounting:development-accounting}
\index{development accounting}\index{levels accounting}
\emph{Development accounting} (or \emph{levels accounting}) splits
differences in income per worker between countries into the part due to
differences in capital and human capital per worker and the part due to
productivity. One summary is the share of the variance of $\ln y$ that $\ln y_{KH}$
accounts for, $\mathrm{Var}(\ln y_{KH})/\mathrm{Var}(\ln y)$.
\end{definition}

\begin{omchart}{Output per worker in 2019 against the output predicted from
capital and human capital alone, $k^{1/3}h^{2/3}$ scaled to the same mean, 144
countries (PPP dollars, log scales), with the 45-degree line.}{Penn World
Table 11.0 (cgdpo, cn, emp, hc)}
\begin{tikzpicture}
\begin{axis}[omscatter, xmode=log, ymode=log, log ticks with fixed point, xtickten={3,4,5}, ytickten={3,4,5},
  xmin=1000, xmax=300000, ymin=1000, ymax=300000, xlabel={Predicted from $K$ and $H$}, ylabel={Actual},
  width=0.56\linewidth, height=0.5\linewidth]
  \addplot[only marks, mark=*, mark size=1.3pt, omDef, fill opacity=0.6] table[x=pred, y=y, col sep=comma] {figdata/bachelor-1/21-growth-accounting/out/devacc.csv};
  \draw[gray, dashed] (axis cs:1000,1000) -- (axis cs:300000,300000);
  \addplot[only marks, mark=none, nodes near coords, point meta=explicit symbolic, nodes near coords style={font=\scriptsize, anchor=west}]
    table[x=pred, y=y, meta=country, col sep=comma] {figdata/bachelor-1/21-growth-accounting/out/devacc_labels.csv};
\end{axis}
\end{tikzpicture}
\end{omchart}

Inputs predict the ranking well --- the correlation of the logarithms is
$0.932$ --- but not the spread: capital and schooling account for $29.6\%$ of the
variance of $\ln y$. % figdata: devacc_stats
The predicted incomes are bunched; the actual ones are not. Relative to the
United States in 2019, India's output per worker was $0.139$ and its inputs
predicted $0.376$, so its productivity was $0.37$ of the American; Niger's
was $0.171$. % figdata: devacc_points
Differences in physical capital and schooling explain only part of income
gaps; most is a difference in the level of the residual (Hall and Jones
1999). % ledger: b1:hall1999

\section{Misallocation}\label{sec:b1:growth-accounting:misallocation}

Aggregate productivity is not only the technology of each firm. It also
depends on which firms use the inputs.

\begin{definition}[Misallocation]\label{def:b1:growth-accounting:misallocation}
\index{misallocation}
There is \emph{misallocation} when the same total inputs, allocated
differently across firms, would produce more; with diminishing returns
within firms, the allocation is efficient only when every firm has the
same \omterm{def:b1:technology:marginal-product}{marginal product} of each input.
\end{definition}

\begin{proposition}[Equal marginal products]\label{prop:b1:growth-accounting:equal-mp}
Let two firms produce $y_i = A_i f(l_i)$ with $f$ increasing and strictly
concave, and $l_1 + l_2 = L$. Total output is maximal when $A_1 f'(l_1) = A_2 f'(l_2)$, and
moving a worker from the firm with the lower \omterm{def:b1:technology:marginal-product}{marginal product} to the one
with the higher raises output.
\end{proposition}

\begin{proof}
Output is $\Phi(l_1) = A_1 f(l_1) + A_2 f(L - l_1)$, with $\Phi'(l_1) = A_1 f'(l_1) - A_2 f'(L - l_1)$ and $\Phi'' < 0$.
The maximum is where $\Phi' = 0$. If $A_1 f'(l_1) < A_2 f'(l_2)$, then $\Phi' < 0$ and lowering $l_1$ raises
output.
\end{proof}

\begin{omfigure}
\begin{tikzpicture}
\begin{axis}[omrecord, grid=none, xlabel={Workers in firm 1 (the rest, of 100, in firm 2)},
  ylabel={Marginal product}, xmin=0, xmax=100, ymin=0, ymax=0.25, width=0.6\linewidth, height=0.4\linewidth,
  samples=80, scaled y ticks=false, yticklabel style={/pgf/number format/fixed}]
  \addplot[demand, domain=4:100, name path=one] {0.5*x^(-0.5)};
  \addplot[supply, domain=0:96, name path=two] {1.0*(100 - x)^(-0.5)};
  \addplot[guide] coordinates {(50,0) (50,0.1414)};
  \addplot[guide] coordinates {(20,0) (20,0.1118)};
  \addplot[omRecord, fill opacity=0.3, draw=none] fill between[of=two and one, soft clip={domain=20:50}];
  \node[font=\footnotesize, anchor=south west] at (axis cs:70,0.06) {firm 1};
  \node[font=\footnotesize, anchor=south east] at (axis cs:93,0.2) {firm 2};
\end{axis}
\end{tikzpicture}

{\small Two firms with $y_1 = l_1^{0.5}$ and $y_2 = 2l_2^{0.5}$ share 100 workers. At 50 each the
\omterm{def:b1:technology:marginal-product}{marginal products} are $0.0707$ and $0.1414$; moving 30 workers to firm 2
equalises them at $0.1118$ and adds the shaded area to output, $5.41\%$.}
\end{omfigure}

The same reasoning, applied to all the manufacturing plants of a country,
gives a measure of what \omterm{def:b1:growth-accounting:misallocation}{misallocation} costs. Reallocating capital and
labour so that \omterm{def:b1:technology:marginal-product}{marginal products} were as dispersed as in the United States
would raise manufacturing productivity by $30$ to $50\%$ in China and by $40$ to
$60\%$ in India (Hsieh and Klenow 2009). % ledger: b1:hsieh2009
Part of the residual of a country is the efficiency with which it
allocates what it has.

\section{Planned economies as data}\label{sec:b1:growth-accounting:planned}

Economies in which the state set investment and allocated inputs are a
test of accounting: they invested heavily, and the question is what the
investment produced.

\begin{record}[The Soviet economy in the accounts]
After controlling for investment and human capital, Soviet growth from
1960 to 1989 was the worst in the world. Its decline from 1950 to 1987 can
be accounted for by a falling \omterm{def:b1:technology:marginal-product}{marginal product} of capital, with
productivity growing at a constant rate: capital piled up and returned
less and less (Easterly and Fischer 1995). % ledger: b1:easterly1995
\end{record}

\begin{omchart}{China: output per worker, capital per worker and total
factor productivity, 1952--2019, TFP from 1956 (1978 = 100, log scale); the dashed line marks
1978.}{Penn World Table 11.0 (rgdpna, rnna, emp, rtfpna)}
\begin{tikzpicture}
\begin{axis}[omtime, ymode=log, log ticks with fixed point, xmin=1952, xmax=2019, ymin=20, ymax=5000,
  xlabel={Year}, ylabel={Index, 1978 = 100}, width=0.72\linewidth, height=0.42\linewidth,
  legend style={font=\footnotesize, at={(0.5,-0.24)}, anchor=north, legend columns=3}]
  \addplot+[thick, no marks] table[x=year, y=y, col sep=comma] {figdata/bachelor-1/21-growth-accounting/out/china.csv};
  \addplot+[thick, no marks] table[x=year, y=k, col sep=comma] {figdata/bachelor-1/21-growth-accounting/out/china.csv};
  \addplot+[thick, no marks] table[x=year, y=tfp, col sep=comma] {figdata/bachelor-1/21-growth-accounting/out/china_tfp.csv};
  \draw[gray, dashed] (axis cs:1978,20) -- (axis cs:1978,5000);
  \legend{Output per worker, Capital per worker, TFP}
\end{axis}
\end{tikzpicture}
\end{omchart}

\begin{record}[China after 1978]
From 1956 to 1978 China's output per worker grew by $2.1\%$ a year, its
capital per worker by $4.18\%$, and its measured productivity fell by $0.62\%$ a
year. From 1978 to 2019 output per worker grew by $7.64\%$ a year, capital per
worker by $8.53\%$ and productivity by $3.18\%$. % figdata: china_growth
The reforms began in agriculture: decollectivisation improved productivity
and accounts for about half of agricultural output growth from 1978 to
1984 (Lin 1992). % ledger: b1:lin1992
\end{record}

Accounting cannot say why the residual is low in one place and high in
another. The deeper causes --- the rules under which people invest,
produce and trade --- are the subject of the chapter on institutions in
the third university year.

\section*{Exercises}\label{sec:b1:growth-accounting:exercises}

\begin{exercise}[$\star$]\label{exo:b1:growth-accounting:residual-ex}
Output grows by $4\%$ a year, capital by $6\%$, labour by $1\%$, human capital
not at all, and $\alpha = 0.35$. Compute the \omterm{def:b1:growth-accounting:tfp}{Solow residual} and its share of growth.
\end{exercise}

\begin{exercise}[$\star$]\label{exo:b1:growth-accounting:read}
From the opening chart's table, Germany's productivity term was $0.56$ of
output growth of $1.72$ (1980--2019), and Japan's $0.3$ of $2.33$
(1970--2019). % figdata: sources
Compute each share of growth due to productivity.
\end{exercise}

\begin{exercise}[$\star$]\label{exo:b1:growth-accounting:steady}
In the \omterm{def:b1:solow-model:steady-state}{steady state} of the Solow model with labour-augmenting progress of
$1.5\%$ a year, predict the growth rate of output per worker and the
contribution of capital deepening to it.
\end{exercise}

\begin{exercise}[$\star$]\label{exo:b1:growth-accounting:china-ex}
From the figures of the China record, which of the two accounting terms,
capital per worker and productivity, changed its growth more between
1956--1978 and 1978--2019? % figdata: china_growth
\end{exercise}

\begin{exercise}[$\star\star$]\label{exo:b1:growth-accounting:levels-ex}
Relative to country A, country B has output per worker $0.25$, capital per
worker $0.2$ and human capital per worker $0.6$. With $\alpha = 1/3$, find the part of
the gap due to inputs, B's relative productivity, and the share of the
log gap that inputs account for.
\end{exercise}

\begin{exercise}[$\star\star$]\label{exo:b1:growth-accounting:two-firms}
In the two-firm example of the figure, find the efficient allocation, the
equalised \omterm{def:b1:technology:marginal-product}{marginal product} and the gain in output.
\end{exercise}

\begin{exercise}[$\star\star$]\label{exo:b1:growth-accounting:hours-ex}
With hours in place of persons, the American residual rises from $0.74$ to
$0.84$ and the German (since 1980) from $0.56$ to $0.98$. % figdata: hours
Explain the sign of each change.
\end{exercise}

\begin{exercise}[$\star\star$]\label{exo:b1:growth-accounting:quality}
Machines improve in quality, but the statistics count them by their
number. Predict the bias in the measured contribution of capital and in
the residual.
\end{exercise}

\begin{exercise}[$\star\star$]\label{exo:b1:growth-accounting:variance}
Capital and schooling account for $29.6\%$ of the variance of $\ln y$ across
144 countries in 2019. % figdata: devacc_stats
If productivity were the same everywhere, what would the variance share be?
What does the number say about the importance of productivity?
\end{exercise}

\begin{exercise}[$\star\star\star$]\label{exo:b1:growth-accounting:tornqvist}
Show that for a Cobb--Douglas technology with constant $\alpha$ the yearly
log-difference decomposition is exact, and that when the share varies
the Törnqvist weights $\tfrac12(\alpha_t + \alpha_{t-1})$ make the error of third order
in the yearly changes (with one year's share it is of second order).
\end{exercise}

\begin{exercise}[$\star\star\star$]\label{exo:b1:growth-accounting:induced}
In the Solow model with labour-augmenting progress at the rate $g$, show
that on the balanced path capital per worker grows at $g$, so that \omterm{def:b1:growth-accounting:growth-accounting}{growth
accounting} attributes $\alpha g$ to capital and only $(1 - \alpha)g$ to the residual,
although all the growth is caused by technology.
\end{exercise}

\begin{exercise}[$\star\star\star$]\label{exo:b1:growth-accounting:refute}
What observation would refute the claim that planned economies grew by
accumulation alone?
\end{exercise}

\section{Problem: Brenland's Growth Review}\label{sec:b1:growth-accounting:problem}

\begin{problem}[{Weekend problem --- accounts and a census}]\label{pb:b1:growth-accounting:brenland}
Brenland's accounts (invented) give real GDP (millions of dalens at 2015
prices), employment (thousands), nominal investment and the deflator
(2015 = 100) for 2000--2025. Build the real capital stock by the perpetual
inventory method: $K_{2000} = 3Y_{2000}$ and $K_{t+1} = 0.95K_t + I_t/P_t \times 100$. Take human capital
as constant. The 2025 census lists 500 firms with their workers, capital
and output (thousands of dalens).

\medskip\noindent\textbf{Part I --- Growth accounting, 2000--2025.} Use $\alpha = 0.35$.
\begin{enumerate}
  \item Compute the average growth rate of real GDP ($Y = 197{,}403$ in 2000 and $292{,}390$ in 2025).
  \item Compute the growth rate of employment (3,760 and 4,366).
  \item Compute the growth rate of output per worker.
  \item Compute $K_{2000}$ and real investment in 2000 (nominal 34,200, deflator 77).
  \item Run the perpetual inventory to 2025: find $K_{2025}$ and the growth rate of capital.
  \item Find the capital--output ratio in 2025 and compare it with 3 in 2000.
  \item Compute the contributions of capital and labour.
  \item Compute the \omterm{def:b1:growth-accounting:tfp}{Solow residual} and its share of output growth.
\end{enumerate}

\medskip\noindent\textbf{Part II --- The residual's uncertainty.}
\begin{enumerate}[resume]
  \item Redo the contributions with $\alpha = 0.30$.
  \item Find the residual and its share.
  \item Why does a lower $\alpha$ raise the residual here?
  \item The starting stock $3Y_{2000}$ is a guess. Which way would a larger $K_{2000}$ move $g_K$?
  \item List two errors that the residual absorbs.
\end{enumerate}

\medskip\noindent\textbf{Part III --- Misallocation in the census.} Assume each
firm produces $y_i = A_i(k_i^{0.35}l_i^{0.65})^{0.85}$, with \omterm{def:b1:technology:returns-to-scale}{returns to scale} of $0.85$.
\begin{enumerate}[resume]
  \item Give the census totals of output, capital and workers.
  \item The \omterm{def:b1:technology:marginal-product}{marginal product} of the input bundle $x_i = k_i^{0.35}l_i^{0.65}$ is $0.85\,y_i/x_i$. Give its tenth and ninetieth percentiles and the standard deviation of its logarithm.
  \item Show that the efficient allocation gives every firm the same share $A_i^{1/0.15}/\sum_j A_j^{1/0.15}$ of both $K$ and $L$.
  \item Compute efficient output and the gain.
  \item Recompute the gain with \omterm{def:b1:technology:returns-to-scale}{returns to scale} of $0.9$.
  \item What share of the inputs would the most productive firm receive?
  \item Why would the gain be overstated if output were measured with error?
\end{enumerate}

\medskip\noindent\textbf{Part IV --- Brenland and Kestria in levels.} Relative to
Brenland, Kestria has output per worker $0.6$, capital per worker $0.5$ and
human capital per worker $0.8$. Take $\alpha = 1/3$.
\begin{enumerate}[resume]
  \item Compute Kestria's relative $y_{KH}$.
  \item Compute its relative productivity.
  \item What share of the log gap do inputs account for?
  \item Compare with the world in 2019.
  \item Which would raise Kestria's income more: closing its capital gap or its productivity gap?
  \item State the result: by how much would equalising \omterm{def:b1:technology:marginal-product}{marginal products} in Brenland's census raise output?
\end{enumerate}
\end{problem}
